I have used data from FAME for this investigation.
FAME is a comprehensive set of financial datapoints from UK firms, with instruments appropriate for this analysis,
in particular number of employees and balance sheet data from which economically interesting variables can be derived.
Fame is provided by Moody's, an international credit agency, and collects its data from UK firms (via survey)\todo{Add citation}.
\par
Its international counterpart, Orbis, offers a similar service for international counterparts.
This analysis could therefore be extended to offer insights comparing the spillover effects of UK firms to their international counterparts.
For this investigation, the UK dataset is comprehensive and the most policy-relevant for this investigation.
\par
Raw data was accessed from FAME in February 2025 by Professor Lars Nesheim at UCL for \citep{Irastorza-Fadrique2026}.
This project stored a a series of Excel tables splitting financial balance sheet variables by their category.
Downloadable indicators under 'a1\_ID' and 'a5\_misc' contain fixed datapoints about each firm; 'a2\_key\_finance',
'a3\_assets', and 'a4\_profits' contain yearly financial data about each firm.
I designed a series of Python scripts to ingest this data and store it in a 'DuckDB' database,
separated into fixed and yearly datapoint tables. 'DuckDB' is an open-source, local, relational (SQL) architecture,
designed for processing bulk queries across columns of variables, making it well-suited for econometric data analysis.